\newgeometry{top=12mm,bottom=15mm,inner=20mm,outer=20mm}
\begin{titlepage}
\begin{spacing}{1.4}
%
\begin{minipage}{0.45\textwidth}
    % \begin{figure}[ht]
    \hspace{-1.5cm}
    \includegraphics[width=\textwidth]{../LOGO_UNISI_ORIZZONTALE_NERO_0.pdf}
    % \end{figure}
\end{minipage}
% \hfil
\hspace{-2.5cm}
\begin{minipage}{0.7\textwidth}
    \flushright
    % \small
    {
    \textit{Dipartimento di Scienze Fisiche, della Terra e dell’Ambiente}\\ \textbf{Dottorato in Fisica Sperimentale}
    }
\end{minipage}
%
% \begin{center}
% {\scshape \Large
% Department of \\Physical Sciences, Earth and Environment
% }

\vspace{1cm}

% \rule{100mm}{0.1mm}
% \rule[5mm]{100mm}{0.4mm}

{
    \large
REPORT ON THE PhD ACTIVITIES CARRIED OUT BY Valeria D'Amante
}



The doctoral student Valeria D'Amante, enrolled in the PhD programme in “Fisica Sperimentale” at the University of Siena, cycle XXXVI, carried out the following activities during the PhD course:

% \vspace{1cm}
\textbf{\textsc{RESEARCH ACTIVITIES}:}\\
% \vspace{22mm}

% \begin{spacing}{2.0}
% {\huge \bf Search for resonant double Higgs production in the bb$\tau\tau$ final state at the CMS experiment.}
% \end{spacing}

% \end{center}

% \vspace{8mm}
% \noindent
% \begin{minipage}[t]{0.5\textwidth}
% \bf Supervisor:\\
% Prof. Maria Agnese Ciocci\\

% %Adjunt Supervisors:\\
% %Dr. Konstantin Androsov\\
% \end{minipage}
% \hfill
% \begin{minipage}[t]{0.5\textwidth}\raggedleft
% \bf Candidate:\\
% Valeria D'Amante\\
% Mat. 106371
% \end{minipage}

% \vfill

% \begin{center}
% \large Academic Year 2023-2024
% \end{center}

\end{spacing}


\end{titlepage}
\restoregeometry
